% Generated draft: requires mathematical and source review.
\documentclass[11pt]{article}
\usepackage[margin=25mm]{geometry}
\usepackage{fontspec}
\setmainfont{Latin Modern Roman}
\usepackage{amsmath,amssymb,mathtools,unicode-math}
\setmathfont{Latin Modern Math}
\usepackage{enumitem,needspace,xurl,hyperref}
\setlength{\parindent}{0pt}
\setlength{\parskip}{5pt}
\setlength{\emergencystretch}{3em}
\newcommand{\N}{\mathbb{N}}
\newcommand{\Z}{\mathbb{Z}}
\newcommand{\Q}{\mathbb{Q}}
\newcommand{\R}{\mathbb{R}}
\newcommand{\C}{\mathbb{C}}
\newcommand{\F}{\mathbb{F}}
\newcommand{\A}{\mathbb{A}}
\newcommand{\PP}{\mathbb{P}}
\newcommand{\E}{\mathbb{E}}
\newcommand{\eps}{\varepsilon}
\newcommand{\dd}{\,\mathrm{d}}
\begin{document}
\section*{UnsolvedMath ID: 20000385}
\textbf{Source problem code:} catalogue entry
\subsection{Definitions and mathematical statement}
The original asks whether a previously defined property ``(A)'' fails for \(\Z/4\Z\)-extensions of function fields when counted by genus. The text defining (A) is absent, so its quantifiers and a precise affirmative/negative assertion cannot be reconstructed. The constant field, base function field, tame/wild convention, connectedness, exact versus bounded genus and automorphism weighting are also unspecified.

The attached aid separates two partial claims. Its geometric layer concerns geometrically connected tame cyclic quartic covers of a genus-zero base, in characteristic different from two. Let \(g\) be the genus of the covering curve, \(h\) the genus of its unique quadratic intermediate curve, and \(b\) the degree of the branch divisor where inertia has order two rather than four. It reports
\[
g=3h+b.
\]
For the corresponding nonempty genus-\(g\) Hurwitz strata with fixed \(h\), it reports codimension \(h\) in the intended full Hurwitz family. The genus-zero base and Hurwitz quotient/stabilizer conventions must be retained; the displayed relation is not a formula for arbitrary base curves or wild characteristic-two covers.

The second layer is an auxiliary positive Euler-product probability law at its discriminant-critical parameter for Kummer tuples. Let \(\mu\) denote that law, let \(Q\) be its proper quadratic-image locus, and let \(E=Q^c\) be the conditioned full \(\Z/4\Z\)-image event. The aid reports
\[
0<\mu(Q)<1,
\qquad\operatorname{Cov}_{\mu}(R_v,R_w\mid E)<0
\quad(v\ne w),
\]
where \(R_v\) indicates ramification at place \(v\). This covariance means
\(\E(R_vR_w\mid E)-\E(R_v\mid E)\E(R_w\mid E)\).
The Kummer tuple space, local weights and explicit covariance formula are not supplied, so this is an attributed statement about that auxiliary law, not a defined arithmetic sampling distribution here.

The aid explicitly does not prove the transfer to the full genus-counted arithmetic family: the modulo-four projection, contribution at infinity, automorphism convention and exact-genus coefficient asymptotics remain unproved. These gaps prevent promoting the auxiliary negative covariance to a conclusion about the original undefined property (A).
\subsection{Short English statement}
Test the missing property (A) for genus-counted cyclic quartic extensions. The attached aid reports a tame Hurwitz-stratum genus relation and negative ramification covariance in a separate conditioned Euler-product model, while explicitly leaving the arithmetic coefficient transfer unproved.
\subsection{Sources}
\begin{itemize}
\item Local UnsolvedMath registry, permanent ID \(20000385\), exact public code \texttt{catalogue entry}; canonical location \texttt{aim-algebraic-number-theory-notes.json}, notes[36].
\item AIM workshop ``Arithmetic statistics over finite fields and function fields,'' Question 7; supplied URL \url{https://aimath.org/pastworkshops/arithstatffieldproblems.pdf}. No proposer, year or definition of (A) is supplied; the literature field is empty. The PDF was not consulted here.
\item Attached research attempt 1, ``A genus-ordering obstruction for cyclic quartic covers,'' without complete bibliographic citations or definitions of its critical Euler-product law.
\end{itemize}
\end{document}
